\section{Additional Theory Analysis}
\label{app:theory-analysis}

\paragraph{Proof of conditional target-slice enrichment}
The result concerns a fixed source pool and a fixed group allocation $B_c$. For $x_i\in T_c$ and $x_j\notin T_c$, attributes with $v_{i,a}=v_{j,a}$ cancel from both score components. With $\mathcal{D}_{ij}=\{a\in\selattrs:v_{i,a}\neq v_{j,a}\}$, the differences in Proposition~\ref{prop:rarity-margin} are
\begin{align}
A_{ij}
&=\sum_{a\in\mathcal{D}_{ij}}
\left[(n_a(v_{i,a}\mid c)+\varepsilon)^{-\alpha}
-(n_a(v_{j,a}\mid c)+\varepsilon)^{-\alpha}\right],\label{eq:rarity-advantage}\\
F_{ij}
&=\sum_{a\in\mathcal{D}_{ij}}
\frac{\hat n_a(v_{i,a})-\hat n_a(v_{j,a})}
{\sum_{v\in\vals_a}\hat n_a(v)}.\label{eq:frequency-disadvantage}
\end{align}
These conditions concern marginal attribute counts. A small joint support $|T_c|$ does not imply $A_{ij}>0$.

\begin{tcolorbox}[
    colback=green!2!white,
    colframe=green!20!white,
    arc=3pt,
    boxrule=1pt,
    drop shadow,
    breakable
]
\begin{proof}[Proof of Proposition~\ref{prop:rarity-margin}]
For each target/non-target pair in group $c$, Eqs.~\ref{eq:rarity-term}--\ref{eq:base-attribute-score} give
\begin{align}
\bar{s}_i-\bar{s}_j
&=\sum_{a\in\selattrs}(R_{i,a}-R_{j,a})\notag\\
&\quad-\lambda\sum_{a\in\selattrs}
\frac{\hat n_a(v_{i,a})-\hat n_a(v_{j,a})}
{\sum_{v\in\vals_a}\hat n_a(v)}\notag\\
&=A_{ij}-\lambda F_{ij}.\label{eq:base-score-gap}
\end{align}
When $F_{ij}\leq0$, the assumed $A_{ij}>0$ yields
\begin{equation*}
\bar{s}_i-\bar{s}_j=A_{ij}-\lambda F_{ij}\geq A_{ij}>0.
\end{equation*}
When $F_{ij}>0$, $\lambda<\lambda_c^\star\leq A_{ij}/F_{ij}$ yields
\begin{equation*}
\bar{s}_i-\bar{s}_j
=F_{ij}\left(\frac{A_{ij}}{F_{ij}}-\lambda\right)>0.
\end{equation*}
At every greedy step $k$, candidates in group $c$ receive the same balance bonus, so
\begin{align*}
\tilde{s}_i^{(k)}-\tilde{s}_j^{(k)}
&=\big(\bar{s}_i+C_c^{(k)}\big)-\big(\bar{s}_j+C_c^{(k)}\big)\\
&=\bar{s}_i-\bar{s}_j>0.
\end{align*}
Thus all target examples precede non-target examples in the group-$c$ queue. The first $B_c$ selections from that queue contain
\begin{align*}
|\dataset_{\mathrm{sel}}^{\mathrm{src}}\cap T_c|
&=\begin{cases}
B_c,& B_c\leq |T_c|,\\
|T_c|,& B_c>|T_c|
\end{cases}\\
&=\min\{B_c,|T_c|\}.
\end{align*}
For a uniform random draw of $B_c$ examples from the same $N_c$-element group, each source example has inclusion probability $B_c/N_c$. Consequently,
\begin{equation*}
\mathbb{E}[|S_{\mathrm{rand}}\cap T_c|]
=\sum_{i\in T_c}\mathbb{P}(i\in S_{\mathrm{rand}})
=B_c\frac{|T_c|}{N_c}.
\end{equation*}
\end{proof}
\end{tcolorbox}

For $0<|T_c|<N_c$ and $0<B_c<N_c$, the conditional enrichment gap is strict:
\begin{equation*}
\min\{B_c,|T_c|\}-B_c\frac{|T_c|}{N_c}
=\begin{cases}
B_c(1-|T_c|/N_c),& B_c\leq|T_c|,\\
|T_c|(1-B_c/N_c),& B_c>|T_c|
\end{cases}>0.
\end{equation*}
This is a selection-count result for a fixed $B_c$. It does not establish that the algorithm assigns positive budget to group $c$, nor does it imply a gain in the trained model's task metric. The random comparison requires the same group allocation and uniform sampling within the group.

A rare conjunction may fail the pairwise condition. For example, if binary combinations $(A,B)=(1,1),(1,0),(0,1),(0,0)$ occur $1,49,49,1$ times, the target conjunction $(1,1)$ has $1\%$ support, but each defining marginal value occurs 50 times. Its marginal rarity score therefore does not strictly separate it from all non-target examples. The sufficient condition $A_{ij}>0$ must be checked for the actual selector fields and score parameters before applying the bound to a particular run.

\paragraph{Geometric coverage}
Let $q_c$ be the empirical $\epsilon$-covering number of $\{z_i:x_i\in T_c\}$ when centers must come from that set. Then $q_c$ target representatives suffice to cover the target slice at radius $\epsilon$. This local fact does not bound the number of target examples in a globally optimal fixed-budget $k$-center solution or in greedy CoreSet: target centers can also cover non-target points, and the global objective can favor additional target centers.

\paragraph{Finite-weight group balance}
Here $c$ denotes a selection group. At step $k$, let $H_c^{(k)}$ be the highest remaining base score in active group $c$ and define $\Delta_k=\max_c H_c^{(k)}-\min_c H_c^{(k)}$. The count $n_c^{(k)}$ is measured before the step.
\begin{tcolorbox}[
    colback=blue!2!white,
    colframe=blue!20!white,
    arc=3pt,
    boxrule=1pt,
    drop shadow,
    breakable
]
\begin{lemma}[Balance pressure at one selection step]
\label{lem:finite-weight-balance}
Assume $0<w<\infty$. If the selector chooses a candidate from group $d$ while group $c$ is also active, then
\begin{equation}
\label{eq:finite-balance-gap}
n_d^{(k)}-n_c^{(k)}\leq\frac{t\Delta_k}{w}.
\end{equation}
\end{lemma}
\end{tcolorbox}

\begin{tcolorbox}[
    colback=green!2!white,
    colframe=green!20!white,
    arc=3pt,
    boxrule=1pt,
    drop shadow,
    breakable
]
\begin{proof}
The chosen group's adjusted head score is at least that of group $c$. Hence
\begin{align*}
H_d^{(k)}+w\frac{t-n_d^{(k)}}{t}
&\geq H_c^{(k)}+w\frac{t-n_c^{(k)}}{t},\\
\frac{w}{t}\big(n_d^{(k)}-n_c^{(k)}\big)
&\leq H_d^{(k)}-H_c^{(k)}\\
&\leq\Delta_k.
\end{align*}
Dividing by $w/t>0$ gives Eq.~\ref{eq:finite-balance-gap}.
\end{proof}
\end{tcolorbox}

The bound applies to counts before the chosen step; it is not a final group-quota guarantee. In the limit $w\to\infty$, finite base scores make the selector choose among active groups with the smallest count. If every group has candidates throughout selection, this gives round-robin group counts differing by at most one. The group bonus is shared within each group, so it preserves that group's base-score ordering.

\subsection{No-Generation Corner Case}
\label{app:no-generation-corner-case}
When $\dataset_{\mathrm{gen}}=\emptyset$, $\tilde n_a^{\mathrm{gen}}(v)=0$ and $\hat n_a(v)=n_a^{\mathrm{src}}(v)$. The source-only score is
\begin{equation}
\label{eq:source-only-base-score}
\bar{s}_i
=\sum_{a\in\selattrs}
\left[
\left(n_a(v_{i,a}\mid g_i)+\varepsilon\right)^{-\alpha}
-\lambda\frac{n_a^{\mathrm{src}}(v_{i,a})}
{\sum_{v\in\vals_a}n_a^{\mathrm{src}}(v)}
\right].
\end{equation}
The second term now penalizes common source attribute values. Proposition~\ref{prop:rarity-margin} still applies if its pairwise score condition holds under these source-only counts. Source selection can only return examples already present in the source pool.
